\documentclass[10pt,a4paper]{amsart}
\usepackage[margin=19mm]{geometry}
\usepackage{amsmath,amssymb,mathtools}
\usepackage[scr=rsfs,cal=euler]{mathalfa}
\usepackage{xfrac}
\usepackage[unicode,hidelinks,bookmarks=false]{hyperref}
\newcommand{\ie}{i.e.\ }
\newcommand{\cf}{cf.\ }
\newcommand{\resp}{respectively\ }
\newcommand{\Subsection}{\subsection}
\providecommand{\nobreakdash}{\nobreak\hbox{-}}
\setlength{\emergencystretch}{2em}
\multlinegap0pt
\usepackage{bm}
\usepackage{bbm}
\usepackage{dsfont}
\usepackage{xspace}
\usepackage{cancel}
\usepackage{mathtools}
\usepackage{amsthm}
\makeatletter
\@ifundefined{theorem}{\newtheorem{theorem}{Theorem}}{}
\@ifundefined{definition}{\newtheorem{definition}[theorem]{Definition}}{}
\@ifundefined{lemma}{\newtheorem{lemma}[theorem]{Lemma}}{}
\makeatother
\newcommand{\CH}{\mathsf{CH}}
\newcommand{\forceP}{\mathbb{P}}
\newcommand{\forceQ}{\mathbb{Q}}
\title[The Consistency of the $\bf{\Sigma^1\_3}$-Separation Property]{The Consistency of the $\bf{\Sigma^1\_3}$-Separation Property}
\author{Stefan Hoffelner}
\date{Source: arXiv:1912.11811v6 (CC BY 4.0)}
\begin{document}
\maketitle

\noindent\emph{Source and licence.} The Consistency of the $\bf{\Sigma^1_3}$-Separation Property. Version arXiv:1912.11811v6 (CC BY 4.0), distributed under the Creative Commons Attribution 4.0 International licence, as recorded in the arXiv licence metadata for this exact version and shown on the abstract page https://arxiv.org/abs/1912.11811v6. Full rights notice: licenses/arxiv-1912.11811-CC-BY-4.0.txt.\par\smallskip

\noindent\textit{Editorial scope.} Counted unit: the Theorem whose source label is \texttt{None} in the article (source0.tex lines 1012--1014, source label \texttt{None}), delivered together with the article's own complete original proof of it (source0.tex lines 1016--1146, one \texttt{proof} environment of 131 lines). No mathematics is written, repaired or re-derived: statement and proof are the article's own text line for line. Required context retained uncounted: none. Every definition, display and label of those lines is retained unchanged. Omitted from this package and not redistributed: 1--1011, 1015--1015, 1147--1264. Nothing inside a retained excerpt is omitted or altered: every retained body is delivered exactly as the article's lines are written, so no omission receipt of any kind accompanies this package. Citations: the retained excerpts carry no citation command at all, so no bibliography entry of the article is transcribed and no reference is redistributed. Reference adaptation: none was needed and none was made; every reference inside the delivered unit resolves inside it, so no unresolved reference or citation is produced. Printed slips kept literal: no printed word, symbol, hypothesis or number of the counted statement or of its proof was altered. Layout and class: the article's own class is \texttt{article} with packages including inputenc, boxedminipage, amsfonts, amsmath, amssymb, graphicx, amsthm, t1enc, subfig, cancel, textcmds; the wrapper is the lane's pinned \texttt{amsart} editorial wrapper at 10pt on A4, which restates the theorem-like environments the retained text needs and adds the article's own macro definitions that the retained text needs (\texttt{\textbackslash CH}, \texttt{\textbackslash forceP}, \texttt{\textbackslash forceQ}), with extra wrapper packages (bm, bbm, dsfont, xspace, cancel, mathtools). The delivered numbering is the unit's own. Assets: no retained span contains a figure environment, a graphics-inclusion command, a PDF-page inclusion, a table or a TikZ picture; no image file is copied into this package, the package declares no asset dependency and no image file is redistributed with it. Licence: version arXiv:1912.11811v6 of this article is distributed under the Creative Commons Attribution 4.0 International licence, as recorded in the arXiv licence metadata for this exact version and shown on the abstract page https://arxiv.org/abs/1912.11811v6. A full rights notice is delivered at licenses/arxiv-1912.11811-CC-BY-4.0.txt. Characters: every retained excerpt is pure ASCII, so the delivered engine, XeLaTeX with its default Latin Modern text font, reports no missing character.
\begin{theorem}
Starting with $L$ as the ground model, one can produce a set-generic extension $L[G]$ which satisfies $\CH$ and in which the ${\Sigma^1_3}$-separation property holds.
\end{theorem}

\begin{proof}


For the the proof to come, we will redefine the notion of legal forcings.
\begin{definition}
A mixed support iteration is called (0-)legal if it is defined as in Definition 4.2, with the only difference that we code (reals that code) quadruples of the form $(x,0,m,k)$ and $(x,1,m,k)$, where $m,k \in \omega$ and $x$ is (the name of) a real into $\vec{S}$.
\end{definition}
Similar to before, we take 0-legal forcings as the base set of forcings and define gradually smaller families of forcings, which we call $n$-legal forcings with respect to a set $E$ and a bookkeeping function $F$.
Assume that $E$ is a finite list of length $n$ of pairwise distinct pairs of natural numbers $(m_l, k_l)$, $l \le n$,
and that $F$ is a bookkeeping function.
Assume that for every $l \le n$ we do have a notion of $l$-legality with respect to $E \upharpoonright l$, let $(m_{n+1}, k_{n+1})$ be a new pair of natural numbers, distinct from the previous ones. Then we say that $\forceP$ is $n+1$-legal with respect to $E \cup \{ (m_{n+1}, k_{n+1} ) \}$ and $F$ if for every $l \le n$, $\forceP $ is $l$-legal with respect to $E \upharpoonright l$ and $F$ and it obeys the following rules:
\begin{enumerate}
\item $\forceP$ never codes a quadruple of the form $(x,1,m_{n+1},k_{n+1})$ into $\vec{S}$.
\item Whenever $\beta < \gamma $, where $\gamma$ is the length of the iteration $\forceP$, is such that there is a $\forceP_{\beta}$-name $\dot{x}$ of a real and an integer $ i\in\{1,2\}$ such that
\[F(\beta)= (\dot{x},m_{n+1},k_{n+1},i)\] and for  $G$ which is $\forceP_{\beta}$-generic over $W$, $W[G]$ thinks that
\begin{align*}
\exists \forceQ (&\forceQ \text{ is } n \text{-legal with respect to } E \,  \land \\&   \forceQ \Vdash x \in A_{m_{n+1}}),
\end{align*}
where $x=\dot{x}^G$, and $y_{\alpha}=\dot{y}_{\alpha+1}^G$.
Then continuing to argue in $W[G]$, we let
\[\forceP(\beta)= \hbox{Code}((x,0,m_{n+1},k_{n+1}),1).\] 

\item
Whenever $\beta < \gamma $ is such that there is a $\forceP_{\beta}$-name $\dot{x}$ of a real and an integer $ i\in\{1,2\}$ such that
\[F(\beta)= (\dot{x},m_{n+1},k_{n+1},i)\] and for $G_{\beta}$ which is $\forceP_{\beta}$-generic over $W$, $W[G_{\beta}]$ thinks that
\begin{align*}
\forall \forceQ_1 (&\forceQ_1 \text{ is } n \text{-legal with respect to } E  \\& \rightarrow \, \lnot(\forceQ_1 \Vdash x \in A_{m_{n+1}}))
\end{align*}
but there is a forcing $\forceQ_2$ such that $W[G_{\beta}]$ thinks that 
\begin{align*}
\forceQ_{2} \text{ is } n \text{-legal with respect to $E$ and } \\   {\forceQ_2} \Vdash x \in A_{k_{n+1}}
\end{align*}


Then continuing to argue in $W[G_{\beta}]$, we force with
\[\forceP(\beta):= \hbox{Code}((x,0,m_{n+1},k_{n+1}),2).\] 




\item If neither 1 nor 2 is true, then either \[ \forceP(\beta)=\hbox{Code}((x,0,m_{n+1},k_{n+1}),2)\] or 
\[ \forceP(\beta)=\hbox{Code}((x,0,m_{n+1},k_{n+1}),1)\] depending on whether $ i\in\{1,2\}$ in $F(\beta)$ was 1 or 2. Otherwise $\forceP$ uses the trivial forcing at that stage. 

\item 
If $F(\beta) = (\dot{x},m,k,i)$ and for every $\forceP_{\beta}$-generic filter $G$, $W[G] \models \forall l \le n+1 ((m_l,k_l) \ne (m,k))$,
then let
\[ \forceP(\beta)=\hbox{Code}((x,0,m,k),i)\] depending on whether $ i\in\{1,2\}$ in $F(\beta)$ was 1 or 2.  
\end{enumerate} 
This ends the definition for the successor step $n \rightarrow n+1$.

With this new notion of $n$-legality, we start the proof of the theorem. The ground model over which we form an iteration is the universe $W$ again, which was defined earlier. Over $W$ we will perform first an $\omega$-length, finitely supported iteration $(\forceP_n)_{n \in \omega}$ of legal posets, and then a second legal iteration of length $\omega_1$. The codes of the form $(x,0,m,k)$ shall eventually define the separating sets for $\varphi_m$ and $\varphi_k$; codes of the form $(x,1,m,k)$ shall correspond to countable sets of reals  (i.e. reals themselves) which indicate the correctness of certain codes of the form $(x,0,m,k)$ which avoid the coding areas coded by these reals.

We let $\{(\varphi_{m_n}, \varphi_{k_n}) \, : \, n \in \omega \}$ be an enumeration of all pairs of $\Sigma^1_3$-formulas. Assume that $(\varphi_{m_0}, \varphi_{k_0})$ is such that there is no legal forcing $\forceQ$ such that $W[\forceQ] \models \exists z (\varphi_{m_0}(z) \land \varphi_{k_0}(z))$. Repeating the arguments from before, we set $E_0:=\{m_0,k_0) \}$ and define the notion of 1-legal with respect to $E_0$.
As will become clear in a second, every step of the iteration $(\forceP_n \, : \, {n \in\omega})$ will either use a legal forcing or define a new and gradually stronger notion of legality. We let $l_n \in \omega$ denote the degree of legality we have already defined at stage $n \in \omega$ of our iteration and define $l_0$ to be $0$ (where 0-legal should just be legal) and $l_1$ to be 1 for the base case of our induction; likewise $\forceP_0$ is set to be the trivial forcing.

The forcing $\forceP_{\omega}$ is the countably supported iteration of $(\forceP_n \, : \,n \in \omega)$, which we will define inductively.
Assume we are at stage $n  \ge 1 \in \omega$ of the iteration and we have defined already the following list of objects and notions:
\begin{enumerate}
\item $\forceP_{n-1}$ and the generic filter $G_{n-1}$.
\item A natural number $l_n \le n$ and a notion of $l_n$-legal relative to $E_{l_n}= \{ (m'_0,k'_0),...,(m'_{l_n-1},k'_{l_n-1}) \} \subset \{ (m_0,k_0),...,(m_{n-1},k_{n-1}) \} $, which is a strengthening of 1-legal relative to $E_0$.
\item A finite set of reals  $\{R_0<...<R_{n - (l_n -1)} \}$, where each real $R_i$ codes a countable set of reals. The choice of the indices will become clear later.

\end{enumerate}
Consider now the $n+1$-th pair $(\varphi_{m_n}, \varphi_{k_n})$ and split into cases.

\begin{enumerate}
\item[$(a)$] There is an $l_n$-legal forcing $\forceQ$ such that \[W[G_{n-1}] \models \forceQ \Vdash \exists z (\varphi_{m_n} (z) \land \varphi_{k_n} (z)), \] then we must use the forcing $\forceQ$. We collect all the reals we have added so far generically which witness $({\ast} {\ast} {\ast})$ for a triple $(x,m,k)$ and call the set $R_{n-l_n}$. In a second step, we use the usual method to code the quadruple $(R_{{n-l_n}}, 1,m'_{l_{n-1}},k'_{l_{n-1}} )$ into $\vec{S}^1$. 
\item[$(b)$] In the second case there is no $l_{n}$-legal forcing $\forceQ$ relative to $E_{l_n}$ which forces $\varphi_{m_n}$ and $\varphi_{k_n}$ to have non-empty intersection. In that case we force with the trivial forcing and define the notion $l_{n+1}$-legal. We first let $(m'_{l_n},k'_{l_n})=(m_n,k_n)$ and $E_{l_n +1}:= E_{l_n} \cup \{(m_n,k_n)\}$, and define $l_{n+1}$-legal relative to $E_{l_n+1}$ just as above. 
We  do not define a new $R_{n-l_n}$.
\end{enumerate}

We let $\forceP_{\omega}$ be the inverse limit of the forcings $\forceP_n$ and consider the universe $W[\forceP_{\omega}]$.
We shall and will assume from now on that in $\forceP_{\omega}$, case $(b)$ is applied infinitely many times. 
\begin{lemma}
For every $n \in \omega$, the tail of the iteration $(\forceP_m \, :\, m \ge n)$ is at least an $l_n-1$-legal iteration relative to $E_{l_n}$ as seen from $W[G_n]$.
\end{lemma}
\begin{proof}
This can easily  be seen if one stares at the definition of the iteration. 
At stage $n$ we either follow case $(a)$ which is an $l_n -1$-legal forcing, as we use the $l_n$-legal $\forceQ$ and additionally code $(R_{{n-l_n}}, 1,m'_{l_{n-1}},k'_{l_{n-1}} )$ which results in an $l_n-1$-legal forcing. Or we apply case $(b)$, thus define $l_n+1$-legality and every further factor of the iteration must be $l_n$-legal. As mixed support iterations of $l_{n}-1$-legal forcings yield an $l_n -1$-legal forcing, this ends the proof.

\end{proof}

The arguments which are about to follow will depend in detail heavily on the actual form of the iteration $\forceP_{\omega}$ which in turn depends on how the enumeration of the $\Sigma^1_3$-formulas behaves. The theorem will of course be true, no matter how $\forceP_{\omega}$ does look like. To facilitate the arguments and notation, however, we assume without loss of generality from now on that the sequence $(\varphi_{m_n}, \varphi_{k_n})$ is so chosen that in the definition of $\forceP_{\omega}$ the cases $(a)$ and $(b)$ are alternating. Thus whenever $n$ is even then $(\varphi_{m_n}, \varphi_{k_n})$ is such that case $(b)$ has to be applied and whenever $n$ is odd then for $(\varphi_{m_n}, \varphi_{k_n})$ case $(a)$ is the one which one should apply. Via changing the order of $(\varphi_{m_n}, \varphi_{k_n})$, this can always be achieved. Consequentially, at even stages $2n$ of the iteration, we define the new notion of $n+1$-legal, while forcing with the trivial forcing; on odd stages $2n+1$, we use the $n+1$-legal forcing to force that $\varphi_{m_{2n+1}}$ and $\varphi_{k_{2n+1}}$ have non-empty intersection and then form the real $R_{2n+1 -n} =R_{n+1}$ and then  code the quadruple $(R_{n+1},1,m_{2n},k_{2n})$ into $\vec{S^1}$.

A consequence of the last lemma (and our assumption on the form of $\forceP_{\omega}$) is that for every even natural number $2n$, there is a final stage in $\forceP_{\omega}$ where we create new codes of the form $(R,1,m_{n},k_{n})$. Indeed, by the definition of $l_{n+1}$-legal, no codes of the form $(R,1,m_{n},k_{n})$ are added by $l_{n+1}$-legal forcings and we have that
\begin{align*}
R_{n}:=  \{ R \, : \,  (R, 1,m_{2n},k_{2n}) \text{ is coded into } \vec{S^1} \}.
\end{align*}
To introduce a useful notion, we say that a real $x$, which is coded into $\vec{S^1}$ or $\vec{S^2}$, has coding area almost disjoint from the real $R$, if $R$ codes an $\omega_1$-sized subset of $\omega$ and the $\omega$-blocks, where $x$ is coded are almost disjoint from the set of ordinals coded by $R$, in that their intersection is countable.

\begin{lemma}
In $W[\forceP_{\omega}]$ for every $n \in \omega \backslash 0$, there is only one real $R$, namely $R_{n}$ which has a code of the form $(R_n,1,m_{2n},k_{2n})$ written into $\omega_1$-many  $\omega$-blocks of elements of $\vec{S^1}$ almost disjoint from $R_{n-1}$.
\end{lemma}
\begin{proof}
This is just a straightforward consequence of the definition of the iteration and of our assumption on the form of $\forceP_{\omega}$.  For $n=1$, note that $R_0$ is the unique ordinal for which $(R,1,m_0,k_0)$ is coded into $\vec{S^1}$. Then the next forcing $\forceP(2)$ is trivial while 2-legal is defined and the forcing $\forceP(3)$ first forces $\varphi_{m_3}$ and $\varphi_{k_3}$ to intersect without creating any code of the form $(R,1,m_0,k_0)$ or $(R,1,m_2,k_2)$, forms $R_1$ and writes $(R_1,1,m_2,k_2)$ into $\vec{S}^1$ with coding area almost disjoint from $R_0$. As all later factors of the iteration are 2-legal, there will be no new codes of the form $(R,1,m_2,k_2)$, hence $R_1$ is the unique real for which $(R_1,1,m_2,k_2)$ is written into $\vec{S^1}$ withwith coding area almost disjoint from $R_0$.

The argument for arbitrary $n$ works exactly the same way with the obvious replacements of letters.

\end{proof}
\begin{lemma}
In $W[\forceP_{\omega}]$, every real  $R_{n}$ is $\Sigma^1_3$-definable.
\end{lemma}
\begin{proof}
This is by induction on $n$. For $n=0$, $R_0$ is the unique real  for which $(R,1,m_0,k_0)$ is coded into $\vec{S^1}$.
This can be written in a $\Sigma^1_3$-way:
\begin{align*}
x= R_0 \Leftrightarrow \exists r \forall M &(|M|=\aleph_0 \land M \text{ is transitive } \land r \in M \land \omega_1^M=(\omega_1^L)^M \rightarrow \\&M \text{ sees with the help of the coded information in } r \text{ that }\\& x \text{ is the unique real such that } (x,1,m_0,k_0) \text{ is coded in} \\& \text{ a block of } \vec{S^1}).
\end{align*}

Now assume that there is a $\Sigma ^1_3$-formula which uniquely defines $R_{n}$, then, by the last lemma, $R_{n+1}$ is the unique real which has a code of the form $(R,1,m_{2(n+1)},k_{2(n+1)})$ written into $\aleph_1$-many $\omega$-blocks of elements of $\vec{S^1}$ almost disjoint from the set of ordinals coded by $R_n$. Let $\psi$ be the $\Sigma^1_3$-formula which defines $R_{n}$, then
\begin{align*}
x= \alpha_{n+1} \Leftrightarrow  &\psi(R_{n}) \text{ and } \\ &\exists r \forall M (|M|=\aleph_0 \land M \text{ is transitive } \land r \in M \land \omega_1^M=(\omega_1^L)^M \rightarrow \\&M \text{ sees with the help of the coded information in } r \text{ that }\\& x \text{ is the unique real such that } (x,1,m_{2(n+1)},k_{2(n+1)}) \\& \text{is coded in} \text{ $\aleph_1$-many blocks of } \vec{S^1} \\& \text{almost disjoint from the ordinals coded by }  R_{n}).
\end{align*}

\end{proof}
The ordinals $R_{n}$ indicate the set of places we need to exclude in order ro obtain correct codes of the form $(x,0,m_{2n},k_{2n})$ which are written into $\vec{S}$. Indeed the iteration $\forceP_{\omega}$, after the stage where we coded the quadruple $(R_{n},1,m_{2n},k_{2n})$ into $\vec{S^1}$ will be $2n-1$-legal, just by the definition of the iteration, which means that the tail of $\forceP_{\omega}$ will never produce a pathological situation.

Finally we can form our desired universe of the $\Sigma^1_3$-separation property. We let $E_{\omega} = \{(m_{2n},k_{2n}) \,: \, n \in \omega\}$ and force with $W[\forceP_{\omega}]$ as our ground model.  We use a countably supported iteration of length $\omega_1$ where we force every quadruple of the form $(x,0,m_{2n},k_{2n})$, $x \in 2^{\omega}$, into either $\vec{S^1}$ or $\vec{S^2}$ with coding area almost disjoint from $ R_{n}$, according to whether case 2, 3 or 4 is true in the definition of $n$-legal. Note that, by assumption, all pairs $(\varphi_{m_n}, \varphi_{k_n})$, $n$ odd, do have a non-empty intersection.  Let $W_1$ denote the universe we obtain this way.

As a consequence we can define in $W_1$ the desired separating sets as follows.
For a pair $(\varphi_{m_{2n}},\varphi_{k_{2n}})$ we let $\psi_{n}$ be the $\Sigma^1_3$-formula which defines $\alpha_{n}$. Now for any real $x$  we let 
\begin{align*}
x \in D_{m_{2n},k_{2n}} \Leftrightarrow \, &\psi(R_{n}) \text{ and } \\& \exists r \forall M (|M|=\aleph_0 \land M \text{ is transitive } \land r \in M \land \omega_1^M=(\omega_1^L)^M  \\&\rightarrow M \text{ sees with the help of the coded information in } r \text{ that }\\& (x,0,m_{2n},k_{2n})  \text{ is coded into $\aleph_1$-many blocks of } \vec{S^1} \\& \text{almost disjoint from the ordinals coded by }  R_{n}).
\end{align*}
and
\begin{align*}
x \notin D_{m_{2n},k_{2n}} \Leftrightarrow \, &\psi(R_{n}) \text{ and } \\& \exists r \forall M (|M|=\aleph_0 \land M \text{ is transitive } \land r \in M \land \omega_1^M=(\omega_1^L)^M  \\&\rightarrow M \text{ sees with the help of the coded information in } r \\& \text{that } (x,0,m_{2n},k_{2n})  \text{ is coded into $\aleph_1$-many blocks of } \vec{S^2} \\&\text{almost disjoint from the ordinals coded by }  R_{n}).
\end{align*}
Both formulas are $\Sigma^1_3$, hence the $\Sigma^1_3$-separation property holds in $W_1$.

\end{proof}

\end{document}
